% Introducción: importancia, origen, objetivos, alcance, limitaciones y
% metodología. Sin anticipar resultados ni conclusiones (NTC 1486).

Una red descentralizada se define por lo que le falta: no hay un servidor, un
controlador ni un punto desde el cual se conozca el estado completo del sistema.
Cada nodo decide con la información que le llega, y esa información es parcial y
llega tarde. Esta propiedad es lo que hace valiosas a estas redes, porque no
dependen de una infraestructura que pueda fallar, y también lo que las hace
difíciles de evaluar: el comportamiento que importa no está en ningún nodo, sino
en la distancia entre lo que cada uno cree y lo que realmente ocurre.

La evaluación de protocolos de red se apoya mayoritariamente en la simulación,
porque desplegar decenas de radios, repetir las condiciones del canal y provocar
fallos a voluntad resulta costoso o directamente imposible. Sin embargo, la
credibilidad de los estudios de simulación es una preocupación documentada
desde hace dos décadas~\cite{Pawlikowski2002Credibility,Kurkowski2005Incredibles},
y se ha mostrado que simuladores distintos pueden dar resultados muy diferentes
para un mismo protocolo~\cite{Cavin2002Accuracy}. A esto se suma un problema de
fondo: una herramienta de propósito general evalúa un modelo del protocolo escrito
para ella y, además, facilita decidir con información global que ningún nodo real
tendría. Un experimento así puede hablar de una red descentralizada sin
respetar la condición que la define.

El origen de este trabajo es un simulador de la red B.A.T.M.A.N. orientado a
operaciones de búsqueda y rescate, disponible en el repositorio del
tutor~\cite{JofsanchezciSimulador}. Al revisarlo con
la orientación del tutor, el foco del proyecto se desplazó hacia una pregunta más
general: \textit{la necesidad de un simulador hecho a la medida para redes
descentralizadas}. La revisión encontró que la herramienta de partida, concebida para otro propósito,
decidía con información global en tres puntos: los Gateway se movían mirando las
posiciones de todos los nodos, los mensajes se entregaban por el camino más corto
del grafo real sin consultar las tablas de enrutamiento, y no circulaba tráfico
de datos. Corregirlos se volvió parte central del trabajo, y la aplicación concreta
pasó a ser una Red de Expansión de Cobertura.

El objetivo de este documento es presentar el diseño, la construcción y la
evaluación de ese simulador. La herramienta ejecuta en cada nodo simulado las
mismas clases de enrutamiento y de detección de fallos que usa el nodo de red, y
sustituye únicamente la capa de transporte por un modelo del medio radio con
distancia y pisos. Se evalúa con un despliegue de cobertura en un edificio de tres
pisos, sometido a distintos tipos de fallo, con Gateway fijos y móviles, y se
compara con líneas base de movilidad centralizada y con un protocolo de inundación.

La metodología combina el desarrollo iterativo guiado por pruebas automatizadas
con un diseño experimental de corridas por semillas, en el que cada cifra del
documento procede de un lote versionado que puede volver a ejecutarse. El trabajo
se limita a un simulador por pasos de tiempo. No incluye un despliegue en \textit{hardware},
no modela colisiones ni la capa de acceso al medio en detalle, y no mide latencia.
Estas limitaciones se tratan de manera explícita a lo largo del documento, porque
condicionan qué se puede afirmar a partir de los resultados.

El documento se organiza así. El Capítulo~\ref{cap:planteamiento} plantea el
problema, los objetivos y el alcance. El Capítulo~\ref{cap:estado-arte} revisa las
herramientas existentes y delimita la brecha. El Capítulo~\ref{cap:metodologia}
describe el método de desarrollo y el diseño experimental. El
Capítulo~\ref{cap:desarrollo} presenta la arquitectura y los mecanismos del
simulador. El Capítulo~\ref{cap:resultados} reporta y discute los resultados, y el
Capítulo~\ref{cap:conclusiones} responde a los objetivos y propone líneas futuras.
Los anexos indican cómo reproducir los resultados.
